\input{preamble}

\begin{document}

\title{Commutative algebra exam}
\maketitle

\section{Hilbert's Nullstellensatz
and Zariski topology}

\begin{exercise}
\label{exercise-maximal-spectrum}
Let $R=\mathbb{R}[t]$ and $X$ be its
maximal spectrum. Prove that $X$ with its
Zariski topology is homeomorphic
to the maximal spectrum of $\mathbb{C}[t]$.
\end{exercise}

\begin{proof}
It's clear that
the maximal spectrums of $\mathbb{R}[t]$
and $\mathbb{C}[t]$ have the same
cardinality. Choose any bijection between
them. We claim this is a homeomorphism
by being a closed map. This follows
since closed sets in any maximal spectrum
are the finite sets (and the
image of a finite set under
a bijection is finite).

To check a closed set $F$ is finite,
suppose $\mathfrak{m} \in F$.
The basis of closed sets is given
by the complements of $U(f)$,
which are defined by
$\mathfrak{m}' \in U(f) \iff
f \not\in \mathfrak{m}'$.
Thus for the closed set $F=U(f)^c$
we have that $\mathfrak{m} \in F
\iff f \in \mathfrak{m}$.
Now since $K[t]$ is a PID,
we may assume $\mathfrak{m}=(g)$,
so $f \in \mathfrak{m}$ becomes
$g$ is a factor of $f$.
But $f$ can have only finitely many factors.

To check any finite set is closed, we
only use the maximal ideals
$t-a$ for any $a \in \mathbb{R}$ or
$\mathbb{C}$.
\end{proof}

\begin{exercise}
\label{exercise-ring-of-plane-curve-is-pid}
Let $P(x,y)$ be an irreducible polynomial
and $R = \frac{\mathbb{C}[x,y]}{(P)}$.
Prove that all ideals in $R$ are
principal, or find a counterexample.
\end{exercise}

\begin{proof}
This is false. Consider
$R=\mathbb{C}[x,y]/(y-x^2)$
and the ideal $(x,y) \subset
\mathbb{C}[x,y]$.
This ideal contains $y^3-x^2$,
so it's an ideal of $R$,
but it is not principal!
Indeed, even after substituting $y^3=x^2$,
there's no way to obtain $y$ from $x$
or vice versa. [Proof can be improved]
\end{proof}

\section{Noetherian rings}

\begin{exercise}
\label{exercise-intersection-submodules-noet}
Let $M$ be an $A$-module, where $A$
is not necessarily Noetherian,
and $N_1,N_2 \subset M$ submodules
which satisfy $N_1 \cap N_2 = 0$.
Assume that the modules $M/N_1$
and $M/N_2$ are Noetherian.
Prove that $M$ is Noetherian.
\end{exercise}

\begin{proof}[Failed proof]
Let $N$ be a submodule (=ideal) of $M$.
Projecting $N$ to $M/N_1$
we obtain a finitely generated module.
Suppose the generators are
$x_i + N_1$, where $x_i \in N$.
The fact that they are generators means
that for any $x \in N$ we have
$x+N_1=\sum a_ix_i + N_1$,
that is, $x-\sum a_ix_i \in N_1$.
The same goes for $N_2$, say,
$x - \sum b_i x_i' \in N_2$ for any
$x \in N$. But then
$\sum a_i x_i - \sum b_i x_i' \in N_1\cap
N_2$, a contradiction (no…)
\end{proof}

\begin{exercise}
\label{exercise-descending-chains-is-field}
Let $A$ be a ring without zero divisors.
Suppose that any descending chain of
ideals in $A$ stabilizes. Prove
that $A$ is a field.
\end{exercise}

\begin{proof}
Suppose that $A$ is not a field
and let's show that there's a descending
chain of ideals which does not stabilize.
Let $a$ be not invertible (since $A$
is not a field).
Then 
$$
(a)\subset (a^2) \subset \subset (a^3)
\subset…
$$
does not stabilize. Indeed,
if $(a^n)=(a^{n+1})$ we obtain
$a^n=ba^{n+1}\iff a^n(1-ba)=0
\implies  1=ba$, since $A$ is a domain,
which is a contradiction because
$a$ is not invertible.
\end{proof}



\section{Irreducible polynomials and Galois
theory}

\begin{exercise}
\label{exercise-rational-functions-irreducible}
Let $k=\mathbb{C}(t)$ be the field of
rational functions on $\mathbb{C}$.
Prove that for any $n$ there exists
an irreducible polynomial
$P(t) \in k[t]$ of degree $n$.
\end{exercise}

\begin{proof}
Let $P(x) \in k[x]=\mathbb{C}(t)[x]$
be $P(x) = x^n-t$.
This polynomial has no root in
$k=\mathbb{C}(t)$,
for if $(f/g)^n=t$ then 
$n(\operatorname{deg}f-
\operatorname{deg}g)=1$,
which is impossible for $n>1$.
But unfortunately this doesn't prove
that it's irreducible: in $\mathbb{R}[x]$,
the polynomial
$(x^2 + 1)(x^2+2)=x^4+3x^2+2$
is reducible and has no roots.

To prove $x^n + t$ is irreducible we shall
use the Gauss lemma. Forget the above
notation and suppose that
$fg = x^n + t$, where
$f,g \in \mathbb{C}(t)[x]$.
We want to clear the denominators in
the coefficients of $f$ and $g$,
so we multiply by the product $D$ of all
such denominators and obtain
$$
f_1g_1 = (x^n + t)D,
$$
where $D \in \mathbb{C}[t]$ and
$f_1,g_1 \in \mathbb{C}[t][x]$.
Now we apply Gauss lemma, that is, the
content of the product is the product
of the contents. Since $x^n +t$ is primitive
(i.e. has content $1$), then the
content of $f_1g_1$ equals the content
of $D$, but $D$ is a constant
(constants here are polynomials in
$\mathbb{C}[t]$).
We obtain
$$
FG = x^n + t
$$
where $F$ and $G$ are the primitive parts
of $f_1$ and $g_1$. This holds up to
multiplication by a complex number.

Now we reduce modulo $t$ and obtain
$$
[F][G] = [x^n],
$$
and conclude that, as classes,
$$
[F]=[cx^m],\qquad 
[G]=[c^{-1}x^{n-m}].
$$
Thus we see that each of $F,G$
lost their constant term when we reduced,
i.e. that their constant terms are divisible
by $t$. This means that $t^2$ divides
the constant term of the product $FG$.
But this is impossible, since the constant
term of $FG$ is $t$. 

The only way
this makes sense is if $m=0$ or $m=n$.
To conclude we note that reduction modulo
$t$ preserves the degree of both $F$, and 
$G$ because their leading coefficients
multiply to $1$. Thus $m=0$ or $m=n$
imply that $F$ or $G$ is a unit, which
in turn implies so for $f$ or $g$.
\end{proof}

\begin{exercise}
\label{exercise-smooth-irreducible-plane}
Prove that for any $d \in  \mathbb{Z}^{>0}$
there exists an irreducible polynomial
$P \in \mathbb{C}[x,y]$ such that
$P(x,y)=0$ defines a smooth submanifold
in $\mathbb{C}^2$.
\end{exercise}

\begin{proof}
$x^n+y$ since $\partial_y =1$.
This is an irreducible polynomial
because $y$ is irreducible
More precisely,
if $x^n+y = fg$, then either of $f,g$
must have no $y$ variable, say $f \in
\mathbb{C}[x]$,
and the other one must be degree 1 in $y$.
Suppose $g = y g' + g''$ for $g',g'' \in
\mathbb{C}[x]$. But then
$x^n+y = fg = yfg' + f g''$,
which means that $f$ cannot have positive
degree in $x$, but $f \in \mathbb{C}[x]$,
so it is a unit.
\end{proof}

\section{Tensor product of R-modules}

\begin{exercise}
\label{exercise-local-flat-implies-flat}
Let $A$ be a ring without zero divisors,
and $M$ an $A$-module
which is flat over all local rings
$A_\mathfrak{p} \supset A$
obtained by localizing $A$
at prime ideals.
Prove that $M$ is flat.
\end{exercise}

\begin{proof}
Let
$$
\xymatrix{
0\ar[r]
&M_1\ar[r]
&M_2\ar[r]
&M_3\ar[r]
&0
}
$$
be exact. Then consider
$$
\xymatrix{
0\ar[r]
&M_1 \otimes_A M\ar[r]^{\varphi}
&M_2 \otimes_A M\ar[r]
&M_3 \otimes_A M\ar[r]
&0.
}
$$
We need to show that the map
$\varphi: M_1 \otimes_A M 
\to M_2 \otimes_A M$ is injective.
So basically the claim is:
if that map is injective
over all localizations $A_{\mathfrak{p}}$
then it is injective over $A$.

So it's basically two steps:
(1) show that
\begin{equation}
\label{equation-localized-varphi}
(\Ker \varphi)_{\mathfrak{p}}
=
\Ker \varphi^{\mathfrak{p}}
\end{equation}
where $\varphi^{\mathfrak{p}}$
is the induced map over $A_{\mathfrak{p}}$
for any $\mathfrak{p} \subset A$
prime. More precisely:
$$
\xymatrix{
0\ar[r]
&(M_1)_{\mathfrak{p}} 
\otimes_{A_{\mathfrak{p}}}M_{\mathfrak{p}}
\ar[r]^{
\varphi^{\mathfrak{p}}}
&
(M_2)_{\mathfrak{p}} 
\otimes_{A_{\mathfrak{p}}}
M_{\mathfrak{p}}\ar[r]
&(M_3)_{\mathfrak{p}}
\otimes_{A_{\mathfrak{p}}}
M_{\mathfrak{p}}\ar[r]
&0
}
$$
where the localizations of each module
are given by tensoring with $A_{\mathfrak{p}}$
over $A$, e.g. $M_{\mathfrak{p}} =
A_{\mathfrak{p}} \otimes_A M$ by definition.
Thus the above sequence is exactly
$$
\xymatrix{
0\ar[r]
&(M_1 \otimes_A M)_{\mathfrak{p}}\ar[r]
^{\varphi^{\mathfrak{p}}}
&(M_2 \otimes_A M)_{\mathfrak{p}}\ar[r]
&(M_3 \otimes_A M)_{\mathfrak{p}}\ar[r]
&0
}
$$
and we see that
Equation \ref{equation-localized-varphi}
follows by definition.
Indeed, both sets are barely
$$
\left\{
\frac{a}{b}x:\varphi(x)=0,
\frac{a}{b}\in A_{\mathfrak{p}}\right\}.
$$


By the exercise's hypothesis, we know that
$\Ker \varphi^{\mathfrak{p}}=0$.

Then, step (2) is:
show that if any $A$-module $K$
is zero as an $A_{\mathfrak{p}}$-module
for any $\mathfrak{p}\subset A$,
then it is zero as an $A$-module.
Here's AI's proof: consider the annihilator
of an element $m \in M$, it's the
set of scalars that kill it, i.e.
$\operatorname{Ann}(m) =
\{a \in A: am=0\}$. It's nonempty by
hypothesis (yes, the fact that $m$ is
zero over $A_{\mathfrak{p}}$ literally
means that there is $a \not\in \mathfrak{p}$
such that $a m =0$).
This just says that $\operatorname{Ann}(m)
\not\subset \mathfrak{p}$ for any
$\mathfrak{p}$.
But that's not true: any ideal is
contained in some maximal (and thus prime)
ideal. The reason for that is Zorn's lemma:
the set of ideals containing
$\operatorname{Ann}(m)$ is inductive
(i.e. every chain is bounded, in this case
by the union of the ideals), so it has
a maximal element, and maximal ideals
are prime because the quotient is a field,
which is in particular a domain.
 
\end{proof}

\begin{definition}
\label{definition-boolean-ring}
A {\it boolean ring} is a ring such
that all its elements are idempotents,
that is, satisfy $a^2=a$.
\end{definition}

\begin{exercise}
\label{exercise-boolean-ring-flat}
Prove that any finitely generated module
over a Boolean ring is flat.
\end{exercise}

\begin{proof}[Failed proof]
The claim is false.
Consider the boolean ring
$\mathbb{Z}/2\mathbb{Z} \oplus
\mathbb{Z}/2\mathbb{Z}$,
equipped with the usual ring structure
for direct sums which is given
by applying the operations entry by entry.

This module has torsion:
$$
([1],[0])([0],[1])=([0],[0]).
$$
Ah! But this ring is not a domain. Forget it.
\end{proof}

\section{Tensor product of rings}

\begin{exercise}
\label{exercise-tensor-product-fg-or-not}
Let $R_1, R_2$ be rings over $\mathbb{C}$
such that $R_1 \otimes_{\mathbb{C}}R_2$
is finitely generated. Prove that
$R_1,R_2$ are finitely generated,
or find a counterexample.
\end{exercise}

\begin{proof}
Suppose 
$R_1 \otimes_{\mathbb{C}}R_2
= \mathbb{C}[z_1,…,z_n]$
and
$z_j=\sum_{k}a_{jk} \otimes b_{jk}$.
Let $A\subset R_1$ be the algebra
generated by the $a_{jk}$.
It's clear that the generators
$z_j$ are in $A\otimes_{\mathbb{C}}R_2$.
Thus, any polynomial on the $z_j$
lies in $A \otimes_{\mathbb{C}}R_2$.
That is, $R_1 \otimes_{\mathbb{C}}R_2
=A \otimes_{\mathbb{C}}R_2$.

Let $\lambda:R_2 \to \mathbb{C}$
be a linear functional with $\lambda(1)=1$.
Consider the operator
$\operatorname{id} \otimes \lambda:
A \otimes R_2 \to A$.

Now take $r_1 \in R_1$.
Then $r_1 \otimes 1 \in 
R_1 \otimes_{\mathbb{C}}R_2
=A \otimes_{\mathbb{C}}R_2$.
Therefore
$r_1=\operatorname{id}\otimes \lambda
(r_1 \otimes 1) \in A$,
and $A$ is finitely generated.
\end{proof}

\begin{exercise}
\label{exercise-trans-deg-hol-functions}
Let $A$ be the ring of complex-analytic
functions on $\mathbb{C}$,
and $k(A)$ its fraction field.
Prove that the transcendence degree
of $A$ over $\mathbb{C}$ is infinite.
\end{exercise}


\section{Normality, irreducibility,
smoothness}

\begin{exercise}
\label{exercise-normal-surface}
Prove that the ring 
$\mathbb{C}[x,y,z]/(y^2-xz)$
is integrally closed.
\end{exercise}

\begin{proof}[Incomplete proof]
Consider a parametrization of the
corresponding affine surface:
$$
(t,s) \longmapsto (t^2,t s,s^2)
$$
This gives an isomorphism
$$
\frac{k[x,y,z]}{(y^2-xz)}
\cong
k[t^2,t s,s^2] \subset k[t,s].
$$
This is integrally closed:
no polynomial in $t^2,st,s^2$ killed
by a polynomial over the fraction
field can be missing.
To prove this we can check which elements
are not in $k[t^2,s t,s^2]$…
which makes no sense -- we have
to search a long the elements
that are roots of monic polynomials
over the fraction field… I'm not sure
about this plan… finding all the
roots of monic polynomials might not
be efficient.
\end{proof}

\begin{exercise}
\label{exercise-quotient-not-smooth}
Let $\zeta_n \in \mathbb{C}$
be a primitive root of unity
of degree $n$.
Define the action of the cyclic group
$G = \mathbb{Z}/\{0\}Z$ on $\mathbb{C}^2$
with coordinates $x,y$
in such a way that the generator
$t$ maps $x$ to $\zeta_nx$ and
$y$ to $\zeta_n^{-1}y$.
Prove that the quotient $\mathbb{C}^2/G$
is singular.
\end{exercise}

\begin{proof}[Ideas]
My ideas:
\begin{itemize}
\item Let $A:= \{f \in k[x,y]:
f(\zeta_nx,\zeta_n^{-1}y)=f(x,y)\}$.
\item We want to prove that $A$ is not
regular. The candidate is obviously the
origin, so we want to localize at the ideal
$(x,y)$, or something like that,
but in the quotient…
\item Since the fibers of the action
are finite sets, the dimension of $A$
should be 2 (the same as $\mathbb{C}^2$).
\end{itemize}

\noindent
Hints used: ``compute invariant ring,
then Jacobian'', and, for invariant ring
computation: ``exponent difference''.

The hint makes us realize that
$A$ is generated by monomials $x^{\alpha}
y^{\beta}$ such that $\alpha-\beta$ is 
a multiple of $n$, because that's the
way we kill the $n$-th roots of unity.
To apply Jacobian criterion we need to 
see $A$ as a finitely generated algebra,
i.e. as a quotient of $k[x,y]$ by
some ideal. The ideal should
take away anything that is not a generator.
This contains $x,x^2,…,x^i$ for $i<n$,
$y,y^2,…,y^j$ for $j<n$ and also
$x^iy^j$ for $i-j<n$.
\end{proof}

\end{document}
